% !TeX program = xelatex
% Generate figures and numeric macros first with python render.py.
\documentclass[UTF8,fontset=fandol,12pt,a4paper]{ctexart}
\usepackage[left=19mm,right=19mm,top=19mm,bottom=20mm,headheight=14pt,headsep=9pt,footskip=23pt]{geometry}
\usepackage{amsmath,graphicx,booktabs,tabularx,array,colortbl,fancyhdr,hyperref}
\usepackage{xcolor}
\definecolor{ink}{HTML}{163342}\definecolor{green}{HTML}{096C65}\definecolor{pale}{HTML}{E6F2EF}\definecolor{muted}{HTML}{4B6570}\definecolor{amber}{HTML}{AA501F}
\hypersetup{unicode=true,colorlinks=true,urlcolor=green,linkcolor=green,pdftitle={Abilene 原八栋园区：按原规划与重新配置的 IT 容量},pdfauthor={IMMC 学生研读资料}}
\input{values.tex}
\setlength{\parindent}{0pt}\setlength{\parskip}{4.3pt}\setlength{\tabcolsep}{4.5pt}
\renewcommand{\arraystretch}{1.25}\raggedbottom
\newcolumntype{Y}{>{\raggedright\arraybackslash}X}
\newcolumntype{L}[1]{>{\raggedright\arraybackslash}p{#1}}
\newcommand{\titleline}[1]{{\sffamily\bfseries\fontsize{22}{29}\selectfont\color{ink}#1\par}\vspace{.7mm}}
\newcommand{\sect}[1]{\par\vspace{1.8mm}{\sffamily\bfseries\fontsize{13}{18}\selectfont\color{green}#1\par}}
\newcommand{\fine}[1]{{\small\color{muted}#1\par}}
\newcommand{\sref}[1]{{\sffamily\small\color{green}[#1]}}
\newcommand{\metric}[2]{\begin{minipage}[t]{.30\linewidth}{\small\sffamily\color{muted}#1\par}{\sffamily\bfseries\fontsize{23}{28}\selectfont\color{green}#2}\;{\scriptsize MW IT}\end{minipage}}
\newcommand{\chart}[1]{\begin{center}\includegraphics[width=.97\linewidth]{figures/#1.png}\end{center}}
\newcommand{\hd}{\rowcolor{pale}}
\pagestyle{fancy}\fancyhf{}\renewcommand{\headrulewidth}{.35pt}
\fancyhead[L]{\footnotesize\sffamily\color{muted}IMMC 2027\quad 问题一\quad /\quad 原规划与反事实求解}
\fancyhead[R]{\footnotesize\sffamily\color{muted}ABILENE, TEXAS}
\fancyfoot[L]{\footnotesize\sffamily\color{muted}学生研读资料\quad ·\quad 2026 年 11 月 22 日}
\fancyfoot[R]{\footnotesize\sffamily\color{muted}\thepage\,/\,4}
\begin{document}\color{ink}
\titleline{按原八栋规划，IT 容量是多少？}
总包 DPR 列出每栋四个数据厅，每厅 25 MW IT，每栋还另列 3 MW IT。开发商 CEO 在访谈中给出的单栋规模约为 103 MW IT。两项工程数据合起来，八栋设计为 $8\times(4\times25+3)=\mathbf{824}$ MW IT，其中数据厅占 800 MW。\sref{1}
\vspace{2mm}\par
\colorbox{pale}{\begin{minipage}{\dimexpr\linewidth-2\fboxsep\relax}\vspace{2mm}\hspace{2mm}\metric{原规划设备}{\planMW}\hfill\metric{仅按三项资源}{\resourceMW}\hfill\metric{Epoch 阶段估计}{824}\vspace{2mm}\end{minipage}}
\sect{让同一套模型检查，这些设备能不能持续运行}
原规划只计常态网电，燃机作备用。若 1,200 MW 并网额定量可全部持续交付，取设计满载设施倍率 $k=1.41$，则供电可支撑 $1{,}200/1.41=851.06$ MW IT。按第 4 页的取水与土地假设，另两项上限分别为 1,266.15 和 962.75 MW IT。资源上限与已配置设备容量取较小值，得到 \textbf{\planMW\ MW IT}。\sref{2–5}
\chart{plan}
\sect{和参照差多少，差在何处}
Epoch 在 2026 年 11 月 1 日估计八栋为 824 MW IT、1,154 MW 设施功率。工程配置复原与其数值相同，两者共用建筑设计信息，零差值只表明口径相合。资源单独计算则高 \textbf{27.06 MW IT（3.28\%）}。\sref{1,3}
\[
 851.06-824
 =\underbrace{(1{,}200-1{,}154)/1.41}_{+32.62}
  +\underbrace{(1{,}154/1.41-824)}_{-5.56}\quad\text{MW IT}.
\]
\fine{第一项来自并网额定量与阶段设施估计量的差别。第二项反映换算倍率和取整，不能据此认定 46 MW 都是限电或闲置。若要在本情景支撑 824 MW IT，至少需 \minGrid\ MW 持续供电（额定量的 96.82\%）和 412 m³/h 到场取水。}
\newpage
\titleline{允许重新配置，求解器会怎么选？}
电网、水量、$k$ 与技术参数沿用原规划对照组。解除八栋设备数量限制，允许重排七宗运营地，并把整宗保留情景中的一块 \patchVal\ ha 候选地用于新增电源或 IT 配套。电源只能使用该候选块。其余用地继续避让，未知的园区地块未加入模型。\sref{5,6}
\sect{求解得到的设备组合}
仓库的混合电源求解器枚举煤、气各 0–2 台，同时优化风光装机、电池逐时调度和恒定 IT 功率。九种火电台数组合中，四种符合地块面积。假定新增常态机组获准投运、燃料充足，最优为 \textbf{两台各 100 MW 的干冷燃气机组}，用地 20 ha。它们是另建的候选机组，原园区的备用燃机仍作备用。\sref{2,7}
\chart{technology}
在此组合下，供电可以支撑 992.91 MW IT。剩给机房与配套的是 $179.61+23.36-20\approx182.96$ ha，按一期用地强度折算，土地先触顶，得到 \textbf{\bestMW\ MW IT}。与原规划设备相比多 \gainPlan\ MW，与纯网电的资源上限相比多 \gainResource\ MW。
\sect{这比原园区实际能增多少？}
连续最优允许重新布置原有土地。若新楼必须按 103 MW 整栋投运，容量向下取整为 \textbf{\wholeMW\ MW IT（九栋）}，实际能否布置仍需总平面检验。另做一组保留八栋的扩建核算，已有宗地不再释放，只允许在 23.36 ha 上加设备。按后六栋的用地增量折算，最优约为 \textbf{\anchoredMW\ MW IT}，只新增一台燃气机组。
\fine{“保留八栋再扩建”使用已知规模估计后六栋的增量用地，专门回答改造问题。重新规划默认余下候选地可接入 IT 配套。若该块只许发电，重配上限回落到 962.75 MW IT。}
\newpage
\titleline{哪些输入会改变答案？}
\sect{供水和持续网电决定设备能否用满}
图中每格都用同一组资源重新求解。左图保留原八栋设备，右图允许重配。取水强度固定为 0.50 L/kWh IT。每格标出的都是 MW IT。
\chart{sensitivity}
若网电只保证额定量的 90\%，到场水量仍为 633 m³/h，原规划可运行约 \textbf{765.96 MW IT}，重配情景约为 \textbf{907.82 MW IT}。若网电全额，但到场水量只有 400 m³/h，两者都落到 \textbf{800 MW IT}。这两条结果分别说明，补建电源可以缓解电力缺口，补水不足时还得改变用水条件。
\sect{单块土地和用地强度}
新增可用块从 0、10、20 ha 扩大到 23.36 ha，最优值依次为 851.06、921.99、962.75、980.75 MW IT。10 ha 和 20 ha 对应一台、两台干冷燃气机组的用地门槛。若每 MW 的 IT 用地强度比一期代理值高 20\%，同一地块下的重配上限降至 861.96 MW IT。这个输入值得优先做宗地核查。把开发商年化设计 PUE 的 1.2、1.3、1.4 暂代满载 $k$，固定八栋仍由设备数量限在 824 MW IT，自由重配分别为 1,034.36、1,000.03、980.75 MW IT。
\sect{只允许风光储时，换一年天气再算}
设计日的最优风光组合把候选块分给约 6.49 MWac 光伏，配 1.62 MW／6.49 MWh 电池。它在 2025 年 1 月 29 日的 24 小时循环中持续增加约 \textbf{0.117 MW IT}。固定同一组合，改用 2024 全年 8,784 小时，首尾荷电量相同，结果为 \textbf{0.154 MW IT}。两个窗口的增益都不到 0.2 MW IT。场址可给光伏的连续面积和夜间储能共同限制了增益。\sref{7,8}
\fine{气象为 Abilene 市域再分析资料。两个时段各自连续校核，跨年长期保障率仍需更多年份。燃煤的高取水系数、燃气恒定可用与设备建设期均为指定的技术情景。}
\newpage
\titleline{哪些是原始输入，哪些仍待核实？}
{\small\begin{tabularx}{\linewidth}{@{}L{18mm}L{50mm}Y@{}}\toprule
\hd 输入&已有资料&本次模型的处理\\\midrule
土地&毛地 445.15 ha。图纸八栋基底合计 36.04 ha。四个可核界区块内，七宗运营地合计约 179.61 ha。&两座站地按 5.85 ha 单列于此计算边界，宗内设施只记一次。完整 1,100 英亩的范围、可用净地和连通性待 GIS 回传。\sref{4–6}\\
单位用地 $a_s$&前两栋租约 94.963 英亩（\phaseHa\ ha），一期约 206 MW IT。&用一期同边界平均值 \slopeVal\ ha/MW IT 代理边际用地。后六栋增量诊断约 0.2284 ha/MW IT，使用了八栋参照，只用于扩建诊断。\sref{1,4,6}\\
电与 $k$&公开并网 1,200 MW。开发商列年化设计 PUE 约 1.2–1.4。&主算例沿用仓库 $k=1.41$。满载实测 $k$、可持续交付量与限电条款仍缺。\sref{2}\\
水与 WUE&市经理称 2025 年约 100 万、2026 年截至 6 月访谈时约 300 万美制加仑/月，含建设和初次充水。&枯水期合同流量、同期全站 WUE 仍缺。基准将 2030 年市级盈余年化并全额给园区，得 633.07 m³/h，现场 WUE 取 0.50。\sref{9}\\\bottomrule
\end{tabularx}}
\fine{Oracle 当前场址页与 Crusoe 均列每栋维护约 50,000 美制加仑/年，Oracle 2025 简报曾列 12,625。维护水量只涵盖部分用途。若要校准 WUE，需要同一时段、同一供水边界的取水量及 IT 电量。\sref{5,9}}
\sect{原始依据与复算入口}
{\small\renewcommand{\arraystretch}{1.14}\begin{tabularx}{\linewidth}{@{}L{8mm}Y@{}}
{[1]}&\href{https://www.dpr.com/projects/crusoe-abilene-data-center}{DPR 总包方的八栋与数据厅参数}，\href{https://www.youtube.com/watch?v=cH7tiy7s17A\&t=48s}{Crusoe CEO 的 CERAWeek 访谈，约 00:48–01:05}。\\
{[2]}&\href{https://www.crusoe.ai/resources/newsroom/crusoe-expands-ai-data-center-campus-in-abilene-to-1-2-gigawatts}{Crusoe 并网与备用说明}，\href{https://media.ffycdn.net/us/crusoe/PL5TuZz5apXB9pVsd3H1.pdf}{Crusoe 2025 影响报告，Abilene 案例（印刷页 16–17）}。\\
{[3]}&\href{https://epoch.ai/graphs/largest-ai-data-centers-by-power-capacity}{Epoch AI 的 IT 容量定义与估计}。阶段设施量取仓库时间线 2026-11-01 行。\\
{[4]}&\href{https://abilenetx.gov/DocumentCenter/View/39456/Lancium-LLC-Abilene-Second-Amended-and-Restated-Tax-Abatement-Agreement-PDF\#page=25}{市税务协议 Exhibit B-1/2，第 25 页}，\href{https://www.nmrk.com/insights/press-releases/newmark-advises-on-3-4-billion-joint-venture-on-one-million-square-foot-texas-data-center}{Newmark 一期融资与规模说明}。\\
{[5]}&\href{https://www.oracle.com/data-centers/abilene/}{Oracle 当前场址页}，\href{https://arrington.house.gov/uploadedfiles/final_oracle_oai_data_center_fact_sheet_092225b.pdf\#page=2}{Oracle 早期简报}，\href{https://www.mortenson.com/projects/abilene-data-center-development}{Mortenson 变电站披露}。\\
{[6]}&\href{https://github.com/RL-RedLeaf/AIDC-test/tree/5374865a8234cd5608c243abc99dfe73bdd18493/data/geographic_case_abilene_20261122/eight_building_land_review}{前次 GIS 面积账与依据}，\href{https://github.com/RL-RedLeaf/AIDC-test/issues/1\#issuecomment-7075327576}{完整宗地的 QGIS 核查请求}。\\
{[7]}&\href{https://github.com/RL-RedLeaf/AIDC-test/tree/5374865a8234cd5608c243abc99dfe73bdd18493/data/geographic_case_abilene_20261122}{仓库案例与混合电源求解器}。对每个整数台数组合求连续容量和周期电池调度。\\
{[8]}&\href{https://open-meteo.com/en/docs/historical-weather-api}{Open-Meteo 历史天气 API / ERA5}。独立气象年、参数表与代码随本报告入库。\\
{[9]}&\href{https://american-audacity.buzzsprout.com/1821492/episodes/19398965-emily-crawford-abilene-city-manager}{市经理 2026-06-25 访谈，40:51 起}，\href{https://www.twdb.texas.gov/waterplanning/rwp/plans/2026/G/RegionG_2026_RWP_VolumeI.pdf\#page=607}{得州 Region G 水规划}，\href{https://www.crusoe.ai/resources/blog/an-inside-look-at-the-abilene-ai-data-center}{Crusoe 用水披露}。
\end{tabularx}}
\end{document}
